\section*{T\_QGR\_Unique: A Bounded-Grammar Uniqueness Theorem for the QM--GR Reconciliation}

\paragraph{Grade.}
This is a finite-carrier bounded-grammar theorem.  It consolidates Lemma 3
(Step 33) and Lemma 4 (Step 34).  The uniqueness grade is precise:
\[
  \text{unique resolution family}
  \quad+\quad
  \text{unique minimal admissible common refinement up to isomorphism}.
\]
It is not an absolute physical uniqueness theorem, not an unrestricted
frame-transfer certificate, and not a claim that the minimal package is the
only admissible refinement.

\subsection*{Definitions}

\paragraph{Bounded grammar \(G^\ast\).}
In this theorem, \(G^\ast\) is the finite grammar consisting of:
\begin{itemize}
\item the declared carrier over modes
  \((d0,d1,d2,d3)\), with \(d0\) density, \(d1\) phase, \(d2\)
  transport, and \(d3\) curvature;
\item endpoint access quotients
  \[
    q_{QM}=(d0,d1,d2),
    \qquad
    q_{GR}=(d0,d2,d3);
  \]
\item the candidate package
  \[
    L=(d0,d1,d2,d3);
  \]
\item the co-sourced field-readout model in which a field \(\psi\) supplies
  \(\mathrm{Born}[\psi]=(\rho,j)\) and \(T[\psi]=(T00,T0i)\);
\item the native promotion-bridge gates
  \(G_{suff},G_{desc},G_{stab},G_{ctrl},G_{nosmuggle},G_{vis},G_{audit},G_{strict}\);
\item F24 layer-multiplicity resolution families:
  \(\texttt{MemoryLayer}\), \(\texttt{HiddenUpstreamRole}\),
  \(\texttt{BridgeMediatedRole}\), \(\texttt{BudgetedRole}\),
  \(\texttt{ScopedRole}\), \(\texttt{CoarsenedRole}\),
  \(\texttt{OutsideRoleScope}\), and \(\texttt{BlockedNonClosure}\);
\item F51 unification as common refinement;
\item the anti-reductionism guardrail: a lawful move may chart, type,
  recognize, construct, simulate, or audit a layer, but may not smuggle one
  endpoint as a derivation of the other across a foreclosed boundary.
\end{itemize}

\paragraph{Access quotient.}
An access quotient is a finite readout map from the carrier to an observable
coordinate package.  In this theorem the endpoint access quotients are
\(q_{QM}\) and \(q_{GR}\).

\paragraph{Common refinement and F51 join.}
A package \(P\) is a common refinement of \(q_{QM}\) and \(q_{GR}\) when both
endpoint quotients factor through \(P\).  In the coordinate-lattice grammar,
the minimal such package is the join of the endpoint-distinguished directions.
For the declared data, this join is
\[
  \{d0,d1,d2\}\cup\{d0,d2,d3\}
  =
  \{d0,d1,d2,d3\}=L.
\]

\paragraph{Admissible.}
A candidate mediator is admissible only if it carries both endpoint descents
and passes the native gates, including the no-smuggling gate excluding
duplicated shared coordinates.

\paragraph{Reconciliation.}
A QM--GR reconciliation in \(G^\ast\) is a lawful package through which the
two endpoint access quotients are jointly mediated without directed reduction
or fused-audit smuggling.  Steps 24--28 give the physical-content form of the
candidate as a co-sourcing field package: one field \(\psi\) sources both the
Born/QM and stress-energy/GR readouts.

\paragraph{Unique up to equivalence.}
Two minimal mediators are equivalent when there is an isomorphism of quotient
coordinates between them preserving the endpoint descents.  In Step 34 this is
exhibited by a permutation \(P\) with residual \(0\) in both directions
between \(L\) and \(L'=PL\).

\subsection*{Theorem}

Within \(G^\ast\), the co-sourcing common-refinement field package \(L\) is
the unique minimal lawful QM--GR reconciliation up to equivalence.

More precisely:
\begin{enumerate}
\item Type uniqueness: the F24 role-split resolves only as
  \(\texttt{BridgeMediatedRole}\) for the actual maps of \(L\).
\item Instance uniqueness: inside \(\texttt{BridgeMediatedRole}\), \(L\) is
  the unique minimal admissible common refinement of \(q_{QM}\) and
  \(q_{GR}\), up to isomorphism.
\end{enumerate}

The theorem does not say that \(L\) is the only admissible object.  Step 34
constructs a five-mode admissible refinement \(M=L+d4\); it is lawful at the
tested gate level but non-minimal because \(L\) factors through it.

\subsection*{Proof}

\paragraph{Lemma 3: type uniqueness.}
Step 33 applies all eight structural F24 predicates to the actual maps of
\(L\).  Exactly one predicate fires:
\[
  \texttt{BridgeMediatedRole}(L)=\top.
\]
The witness residual for the explicit joint quotient is \(0.0\), and the
native gate proxy has pass count \(8\).  The remaining seven F24 families are
computed-excluded.

The Step 33 coarsening search checks all \(8\) quotients obtained by retaining
subsets of \(q_{QM}\)'s modes.  Every such coarsening leaves \(O_s\) nonempty;
the minimum obstruction count is \(3\).  Thus
\(\texttt{CoarsenedRole}\), the finite analogue of a coarsened-only
reconciliation, is excluded as the full reconciliation type in this grammar.
The other exclusions are likewise structural: no hidden latent is present,
the resolution is exact rather than residual-only, no proper scope restriction
is used, no history coordinate carries the split, \(d3\) is in scope, and a
closure forms.

Sources:
\texttt{steps/step33\_partb\_type\_uniqueness\_artifacts/step33\_results\_summary.md},
\texttt{steps/step33\_partb\_type\_uniqueness\_artifacts/f24\_predicates\_on\_L\_step33.csv},
and
\texttt{steps/step33\_partb\_type\_uniqueness\_artifacts/excluded\_family\_rulings\_step33.csv}.

\paragraph{Lemma 4: instance uniqueness.}
Step 34 computes the common-refinement lattice around \(q_{QM}\) and
\(q_{GR}\).  Both endpoints factor through \(L\) with residual \(0.0\), and
the directions of \(L\) are exactly the union of the endpoint directions.  So
\(L\) is the coordinate-lattice join.

Every drop-one-mode subobject fails to carry both endpoints by positive linear
obstruction.  The missing-mode endpoint residual in each failure case is
\[
  0.5773502691896258.
\]
The five-mode refinement \(M=L+d4\) is admissible but non-minimal:
\[
  \mathrm{res}(L\leftarrow M)=0.0,
  \qquad
  \mathrm{res}(M\leftarrow L)=0.4472135954999579.
\]
The direct-sum union
\[
  (d0,d1,d2,d0,d2,d3)
\]
is not admissible: it has rank \(4\) in dimension \(6\), with extra coordinate
count \(2\).  Finally, a relabeled copy \(L'=PL\) is isomorphic to \(L\) by an
exhibited permutation, with residual \(0.0\) in both directions.

Sources:
\texttt{steps/step34\_partb\_instance\_uniqueness\_artifacts/results\_summary.md},
\texttt{steps/step34\_partb\_instance\_uniqueness\_artifacts/mediator\_candidates\_step34.csv},
\texttt{steps/step34\_partb\_instance\_uniqueness\_artifacts/descent\_residuals\_step34.csv},
and
\texttt{steps/step34\_partb\_instance\_uniqueness\_artifacts/isomorphism\_check\_step34.csv}.

\paragraph{Controls.}
Step 33's predicate computation is not tag-driven: the family firing is
computed from \(L\)'s actual maps, and a second firing family would be reported
as non-uniqueness.  Step 34's super-minimal control shows that the uniqueness
claim is not vacuous: other admissible mediators exist, but they are
refinements above the minimal join rather than second minimal mediators.  The
drop-one-mode controls show that proper subobjects fail to carry both endpoint
quotients.

\paragraph{Conclusion.}
Lemma 3 gives unique resolution type.  Lemma 4 gives unique minimal instance
within that type, up to isomorphism.  Therefore \(L\) is the unique minimal
lawful QM--GR reconciliation in \(G^\ast\), up to equivalence.

\subsection*{Relation to Part A}

T\_QG\_NoGo (Step 32) proves that the directed and fused single-object routes
are not lawful in \(G^\ast\).  T\_QGR\_Unique proves that the lawful
co-sourcing common-refinement route is unique at the bounded minimal
up-to-equivalence grade.  Together, the two Mode-T theorems complete the
bounded program: the fused/derived target is foreclosed in \(G^\ast\), and the
co-sourcing common-refinement is forced as the unique minimal lawful
reconciliation in \(G^\ast\).

\subsection*{Seven-Gate Audit}

\begin{enumerate}
\item Self-contained definitions: \(G^\ast\), access quotient, common
  refinement, F51 join, admissibility, reconciliation, and equivalence are
  defined above.
\item Lemma evidence: Lemma 3 cites Step 33 computation; Lemma 4 cites Step 34
  computation.
\item Controls: Step 33 structural predicate controls and Step 34
  sub-minimal, super-minimal, union, and isomorphism controls are cited.
\item Bounded grade: finite-carrier bounded-grammar uniqueness, up to
  equivalence, at the unique-minimal common-refinement grade.
\item No-smuggling: the direct-sum union's \(G_{nosmuggle}\) failure is part of
  the instance proof.
\item Scope boundary: no statement beyond \(G^\ast\) is asserted.
\item Next obligation: external frame-transfer review
  \(G\_E018\_FRAME\_TRANSFER\_EXTERNAL\_REVIEW\_STEP35\).
\end{enumerate}

\subsection*{Interpretation}

Framework-side, the result explains why the co-sourcing semiclassical pattern
is the lawful reconciliation recovered by the construction: it is the minimal
common refinement selected by the grammar.  The excluded F24 families line up
with alternative full-reconciliation programs in the finite analogue: hidden
latent, coarsened-only, residual-only, regime-scoped, memory-only,
outside-scope, and blocked-nonclosure resolutions.  They may remain useful as
local methods or approximations, but they are not the full reconciliation type
for this \(L\) in \(G^\ast\).  See
\texttt{interpretation/quantum\_gravity\_is\_illegal.md}
and
\texttt{interpretation/ladder\_vs\_fork.md}.

\subsection*{Nonclaim}

This theorem is a bounded statement internal to \(G^\ast\).  It does not make
an absolute physical uniqueness claim, does not assert a complete physical
model, does not derive one endpoint from the other, and does not certify frame
transfer.  The uniqueness notion is exactly: unique resolution family plus
unique minimal admissible common refinement up to isomorphism on the declared
finite carrier.  Super-minimal admissible refinements exist and are
non-minimal.  External review remains the standing next gate.
